\scanpage{48}
% archive foliation: 47
loc. isomorphe à $\{0,\infty,1,-1,i,-i\}$ de $\mathbb P^1_S$, ou ce qui revient au $=$, qui soit loc. isom. à la configuration octaédrale associée à une op. \inserted{monom.} d'un $\underline V$ sur $X$. Nous allons voir maintenant que cette donnée est essentiellement unique.

Pour ceci, pour un diviseur octaédral \deleted{$\unsure{\Delta}$}
\[
\Delta\subset X
\]
considérons le schéma en groupes
\[
(7)\qquad\Gamma=\underline{\operatorname{Aut}}_S(X,\Delta),
\]
Il est immédiat (du fait que $\Delta$ est de degré $\geq3$) que $\Gamma$ \inserted{$\subset\underline{\mathfrak S}_{\Delta}=\underline{\operatorname{Aut}}_S(\Delta)$} est \uline{net} sur $S$, Je dis que $\Gamma$ est étale sur $S$, et loc. isom. au schéma en groupes $(\mathfrak S_4)_S$. Plus précisément, je vais définir sur $\Delta$ \inserted{(comme schéma de revêtements)} une structure d'\deleted{\unsure{oct.}} schéma en octaèdres combinatoires, de telle façon que l'on ait
\[
(8)\qquad\Gamma=\underline{\operatorname{Aut}}^+_{\mathrm{oct.}}(\Delta)
\qquad(\text{en tant que sous-schéma de }\underline{\operatorname{Aut}}_S(\Delta))
\]
Pour ceci, on \deleted{\unsure{se donne une}} \inserted{note qu'une} structure octaédrale sur un ens. à \unsure{six|dix} éléments est bien déterminée quand on connaît le groupe de ses automorphismes \inserted{directs} (en tant que sous-groupe de $\mathfrak S_\Delta$), car
